% The three figures the Module 7 sheet draws, kept in their own file because
% the answer key draws exactly the same figures with the answers on them.
% Needs tikz and the colours pltblue and pltrust, which the including file sets
% up.
% ---------------------------------------------------------------------------
\tikzset{
  friend/.style  = {line width=1.3pt, black!58, line cap=round},
  person/.style  = {circle, draw=black!72, fill=white, line width=0.9pt,
                    minimum size=0.74cm, inner sep=0pt},
  % The walker's start: a second ring round the person, as on slide 6.
  start/.style   = {person, double=white, double distance=1.6pt},
  % A box the student writes a number into (a share, in Question 6).
  sbox/.style    = {draw=pltblue!70, line width=0.7pt, fill=white,
                    minimum width=1.25cm, minimum height=0.62cm, inner sep=0pt,
                    rounded corners=1.2pt},
  member/.style  = {circle, draw=black!72, fill=white, line width=0.9pt,
                    minimum size=0.66cm, inner sep=0pt, font=\small},
  dot/.style     = {circle, fill=pltblue!80!black, inner sep=0pt,
                    minimum size=0.22cm},
}

% ---------------------------------------------------------------------------
% Part 1: four friends. A--B, A--C, B--C, A--D. Degrees 3, 2, 2, 1, so 2m = 8
% and the long-run shares are 3/8, 2/8, 2/8, 1/8. D hangs off A, so the walker
% swings between A and the rest; the share column of Question 6 does not care.
%
% #1 is drawn last, on top: the answer key writes the shares into the boxes
% with \node at (bA) {...}; and so on.
\newcommand{\fournet}[1]{%
\begin{tikzpicture}
  \node[start]  (A) at (0, 0)      {A};
  \node[person] (B) at (1.9, 1.25) {B};
  \node[person] (C) at (2.6, -0.45){C};
  \node[person] (D) at (-1.9, 0)   {D};
  \draw[friend] (A) -- (B);
  \draw[friend] (A) -- (C);
  \draw[friend] (B) -- (C);
  \draw[friend] (A) -- (D);
  \node[font=\footnotesize, black!65, above=0.08cm of A] {start};
  \node[sbox] (bA) at (0, -1.05)   {};
  \node[sbox] (bB) at (3.35, 1.25) {};
  \node[sbox] (bC) at (2.6, -1.45) {};
  \node[sbox] (bD) at (-1.9, -1.0) {};
  #1
\end{tikzpicture}}

% ---------------------------------------------------------------------------
% Part 3: twelve members, the network of slides 14-45. Members 0-5 and 6-11
% are two groups of six, everybody friends with everybody in the group except
% 0-3 and 7-10, and two friendships cross: 2-8 and 5-11. 30 friendships.
% The layout is the deck's (slides/m07/figures/smallgraphs.py, TWO12_XY): two
% hexagons, the missing friendships are long diagonals, and the two crossing
% friendships are parallel horizontal lines that touch no other member.
%
% #1 is the style of members 0-5 and #2 of members 6-11 (empty on the sheet:
% the student darkens them; the key fills one side).
\newcommand{\tworooms}[2]{%
\begin{tikzpicture}[x=0.8cm, y=0.8cm]
  \foreach \n/\a in {2/60, 1/120, 0/180, 4/240, 5/300, 3/0}
    \coordinate (m\n) at ($(1.5, 2.0) + (\a:1.7)$);
  \foreach \n/\a in {6/60, 8/120, 10/180, 11/240, 9/300, 7/0}
    \coordinate (m\n) at ($(7.5, 2.0) + (\a:1.7)$);
  \foreach \i/\j in {0/1, 0/2, 0/4, 0/5, 1/2, 1/3, 1/4, 1/5, 2/3, 2/4, 2/5,
                     3/4, 3/5, 4/5,
                     6/7, 6/8, 6/9, 6/10, 6/11, 7/8, 7/9, 7/11, 8/9, 8/10,
                     8/11, 9/10, 9/11, 10/11,
                     2/8, 5/11}
    \draw[friend] (m\i) -- (m\j);
  \foreach \n in {0,...,5}  \node[member, #1] at (m\n) {\n};
  \foreach \n in {6,...,11} \node[member, #2] at (m\n) {\n};
\end{tikzpicture}}

% ---------------------------------------------------------------------------
% Part 3: the empty axes the second eigenvector is plotted on, one column per
% member, -1 to 1. #1 is drawn on top (the key's dots).
\newcommand{\eigaxes}[1]{%
\begin{tikzpicture}[x=1.15cm, y=2.9cm]
  \foreach \y in {-1, -0.5, 0.5, 1}
    \draw[black!14, line width=0.5pt] (-0.5, \y) -- (11.5, \y);
  \foreach \x in {0,...,11}
    \draw[black!14, line width=0.5pt] (\x, -1.12) -- (\x, 1.12);
  \draw[black!70, line width=0.8pt, dashed] (-0.5, 0) -- (11.5, 0);
  \draw[black!80, line width=0.9pt] (-0.5, -1.12) -- (-0.5, 1.12);
  \foreach \y/\l in {-1/{$-$1}, -0.5/{$-$0.5}, 0/{0}, 0.5/{0.5}, 1/{1}}
    \node[left, font=\small] at (-0.5, \y) {\l};
  \foreach \x in {0,...,11}
    \node[anchor=base, font=\small] at (\x, -1.27) {\x};
  \node[font=\small, black!65] at (5.5, -1.5) {member};
  #1
\end{tikzpicture}}
